\chapter{浮动基座降阶模型综览——LIPM·ZMP·DCM·SLIP·SRBD·Centroidal 的模型谱系与选型}
\label{ch:reduced-models}

% ════════════════════════════════════════════════════════════════
%  控制部 C3 第 M 章 · 卷五《机器人最优控制》腿足/接触子部
%  定位：腿足降阶模型的"统一谱系 + 选型理论"章 + 三块独占拼图(SLIP/ZMP预览/VMC) + CMM 桥。
%  源：Tutorial 51(腿足简化模型理论, 主体) + 49.9-49.10(浮动基座/质心动量, 桥)。
%  硬边界(去重)：SRBD-MPC 全栈→\cref{ch:quad_mpc}；DCM 步态/ALIP→\cref{ch:legged-biped}；
%               RNEA/全身动力学→\cref{ch:spatial-dynamics}；CCRBA→\cref{ch:contact-mechanics}。
% ════════════════════════════════════════════════════════════════

\begin{practice}[前置自测]
开始之前，先试答下列问题。若已能流畅作答，可只速读本章的谱系图与选型表；若卡住，正是本章要为你解决的。
\begin{enumerate}
  \item 一台 Go2 四足机器人的全身状态有多少维？若直接对全身动力学做 $N=20$ 步的模型预测控制（MPC），决策变量大约多少？为何难以在毫秒级控制周期内实时求解？
  \item 线性倒立摆（LIPM）方程 $\ddot{x}=\omega_n^2(x-p)$ 为什么描述一个\emph{本征不稳定}的系统？它与“在指尖上立一根棍子”是同一回事吗？
  \item 零力矩点（ZMP）、压力中心（CoP）、质心（CoM）三者有何区别？“ZMP 落在支撑多边形内”是稳定行走的充分条件还是必要条件？
  \item 发散分量（DCM）$\xi=x+\dot{x}/\omega_n$ 把 LIPM 的二阶系统降成了几阶？这个降阶的“本质洞察”是什么？
  \item 单刚体模型（SRBD）相对全身动力学“丢掉”了哪些物理？为什么它的 MPC 能写成凸二次规划，而质心动力学（Centroidal）模型的 MPC 一般只能是非线性规划？
\end{enumerate}
\end{practice}

\section*{本章目标}
\addcontentsline{toc}{section}{本章目标}
学完本章，你应当能够：
\begin{enumerate}
  \item \textbf{说清}全身动力学为何无法直接用于实时 MPC，以及降阶模型通过哪些物理假设把数十维系统压到 $2$–$12$ 维；
  \item \textbf{把握}一张统一的\emph{降阶模型谱系}（LIPM/DCM/SLIP/SRBD/LL-SRBM/Centroidal），理解每个模型\emph{抛弃了什么物理}、何时该升档/降档；
  \item \textbf{推导} LIPM 方程、ZMP 完整公式（含角动量项）与 ZMP$=$CoP 定理，并解释 ZMP 判据“必要非充分”的含义；
  \item \textbf{重建} Kajita 的 ZMP 预览控制（Cart-Table $+$ LQR $+$ preview），说明预览增益为何随预览步数指数衰减；
  \item \textbf{理解} DCM 作为“LIPM 一阶降阶”的本质，以及 SLIP 模型如何用弹簧腿与飞行相突破 LIPM 的恒高假设；
  \item \textbf{掌握} Raibert 三分解耦与落脚启发式，并看清它与捕获点的深层联系；
  \item \textbf{运用}质心动量矩阵（CMM）$\mathbf{h}_G=\mathbf{A}_G(q)\dot{q}$ 作为“全身动力学 $\leftrightarrow$ 简化模型”的精确投影桥，把整张谱系收束到一条主线；
  \item \textbf{使用}虚拟模型控制（VMC）的虚功原理力映射 $\boldsymbol{\tau}=\mathbf{J}^\top\mathbf{F}$，并说明它为何在高动态下不及全身控制（WBC）。
\end{enumerate}

\subsection*{本章在全书中的位置}
\addcontentsline{toc}{subsection}{本章在全书中的位置}
\cref{ch:spatial-dynamics}（空间向量代数 $\cdot$ RNEA $\cdot$ ABA）已经给出浮动基座机器人的\emph{全身}动力学方程 $\mathbf{M}(q)\ddot{q}+\mathbf{h}=\mathbf{S}^\top\boldsymbol{\tau}+\mathbf{J}_c^\top\boldsymbol{\lambda}$；\cref{ch:contact-mechanics}（接触力学）则把接触约束、摩擦锥与质心动量矩阵的 $O(N)$ 算法做透。本章站在这两章之上，回答一个工程上绕不开的问题：\textbf{既然全身动力学维度太高、无法实时优化，规划层该用什么“低维替身”？} 答案是一族\emph{降阶模型}。它们不是某一个孤立模型，而是一张\emph{谱系}——从最简的二维线性倒立摆，到最精确的六维质心动力学，每一档都用一组物理假设换取一截维度与算力。

\begin{note}
\textbf{本章定位（去重声明）.}\;
本章是腿足降阶模型的\textbf{统一谱系与选型理论}章，独占内容为：\emph{谱系总图 $+$ 选型理论}、\emph{SLIP 弹簧倒立摆}、\emph{ZMP 预览控制}、\emph{虚拟模型控制 VMC}、以及\emph{质心动量桥}。各模型的工程深挖见 P7 腿足部对应章——\textbf{单刚体 MPC 全栈}（SRBD 完整推导、三大凸化近似、$13$ 维状态、堆叠 QP、Di Carlo 工程细节）见 \cref{ch:quad_mpc}；\textbf{DCM 步态规划与 ALIP 横向落脚律}见 \cref{ch:legged-biped}；\textbf{空间向量/RNEA/全身动力学}见 \cref{ch:spatial-dynamics}；\textbf{CCRBA 算法/CCRBI}见 \cref{ch:contact-mechanics}。凡触及这些主体，本章只给指针、不重复推导。
\end{note}

\begin{note}
\textbf{本章记号与约定.}\;
本章高频出现两个易混的 $\omega$：LIPM 的\textbf{自然频率}（标量）记 $\omega_n=\sqrt{g/h}$，机体\textbf{角速度}（三维向量）记 $\boldsymbol{\omega}$。发散分量记 $\xi$（平面量，\textbf{专指 DCM，与 \cref{ch:lie} 的 $\mathfrak{se}(3)$ 旋量 $\boldsymbol{\xi}=[\boldsymbol{\rho};\boldsymbol{\phi}]$ 无关}）；压力中心、零力矩点记 $\mathbf{p}_{\mathrm{CoP}}$、$\mathbf{p}_{\mathrm{ZMP}}$，质心记 $\mathbf{r}_G$，第 $i$ 个接触点位置记 $\mathbf{p}_i$、其接触力记 $\mathbf{f}_i$。质心动量统一写成\textbf{线动量在前、角动量在后}：$\mathbf{h}_G=[\mathbf{l}_G;\,\mathbf{k}_G]$，与本书空间力序 $\mathbf{w}=[\mathbf{f};\mathbf{n}]$、空间速度序 $\boldsymbol{\varpi}=[\boldsymbol{\nu};\boldsymbol{\omega}]$ 一致（外部文献 Featherstone/Orin 多用角动量在前的 $[\mathbf{k};\mathbf{l}]$，两者差一个 $6\times 6$ 置换；Pinocchio 的 \texttt{Force} 恰与本书同序）。姿态扰动一律用\textbf{右扰动} $\mathbf{R}\,\mathrm{Exp}(\delta\boldsymbol{\phi})$（\cref{ch:lie}）。
\end{note}

% ════════════════════════════════════════════════════════════════
\section{维度灾难与降阶哲学}
\label{sec:rm-curse}

\paragraph{动机：全身 MPC 为何不实时。}
考虑 Go2 四足机器人的全身状态：广义坐标 $q\in\mathbb{R}^{19}$（$7$ 维浮动基座位姿 $+$ $12$ 维关节角），广义速度 $\dot{q}\in\mathbb{R}^{18}$，控制输入是 $12$ 个电机力矩 $\boldsymbol{\tau}\in\mathbb{R}^{12}$，接触力 $\boldsymbol{\lambda}\in\mathbb{R}^{12}$（$4$ 脚 $\times$ $3$D）。MPC 要在每个控制周期同时优化未来一段时域内的所有这些量。取时域 $N=20$，仅决策变量就有
\begin{equation}
  N\times(n_v+n_\tau+n_\lambda)=20\times(18+12+12)=840,
  \label{eq:rm-decision-count}
\end{equation}
对应一个 $840\times 840$ 的 Hessian。这样规模的二次规划在嵌入式平台上的单次求解时间会远超腿足 MPC 通常的毫秒级周期预算；人形机器人（关节更多）则彻底不可行。\cref{tab:rm-dim} 给出几个平台的数量级对照。

\begin{table}[H]\centering\small
\caption{全身动力学与各降阶模型的维度／算力数量级对照（依实现、平台与求解器而变，仅作直觉）}
\label{tab:rm-dim}
\begin{tabular}{lccccc}
\toprule
模型 & 状态维度 & 决策变量($N{=}20$) & 优化结构 & 实时性 & 典型用途 \\
\midrule
全身动力学（Go2, $18$-DOF） & $37$ & $\sim\!840$ & 大规模 QP/NLP & 一般离线 & 真值参照 \\
全身动力学（H1, $25$-DOF） & $51$ & $\sim\!1200$ & 大规模 QP/NLP & 不可行 & — \\
LIPM & $2$–$4$ & $\sim\!80$ & 极简 LQR/QP & 极高 & 双足步行 \\
SRBD & $12$–$13$ & $\sim\!260$ & 凸 QP & 高 & 四足 trot \\
Centroidal & $12$ & $\sim\!240$ & NLP（SQP/DDP） & 中 & 通用腿足 \\
\bottomrule
\end{tabular}
\end{table}

\paragraph{核心思想：只建模质心，不建模关节。}
降阶的出发点是一个朴素观察：步态规划与平衡控制本质上只关心系统的\emph{整体动量}——质心在哪里、往哪里动、整机转得多快——而不关心每个关节的具体角度。关节角度是“如何实现这个动量目标”的执行细节，可以交给下层处理。于是腿足控制天然分成两级：
\begin{itemize}
  \item \textbf{规划层}（本章的降阶模型 MPC）：决定“质心应从这里移到那里、整机该施加多大合力与合力矩”；
  \item \textbf{控制层}（全身控制 WBC，\cref{ch:control} 与 P7 的 WBC 章）：把规划层给出的“期望质心力/接触力”分解为“各关节力矩”。
\end{itemize}
这一分层把一个数十维的实时优化，拆成“低维 MPC（数十毫秒一拍）$+$ 高频 WBC（千赫兹一拍）”，是现代腿足控制栈的骨架（\cref{fig:rm-pyramid}）。

\begin{figure}[H]
\centering
\begin{tikzpicture}[
  >=Stealth, node distance=3mm,
  lvl/.style={draw, rounded corners, align=center, font=\small, inner sep=5pt, text width=86mm},
  plan/.style={lvl, fill=main!10, draw=main!60!black},
  ctrl/.style={lvl, fill=second!10, draw=second!60!black},
  motor/.style={lvl, fill=third!10, draw=third!60!black},
  ln/.style={->, thick, gray!55!black}]
\node[plan] (plan) {\textbf{规划层：降阶模型 MPC}（$\sim\!30$–$50$\,Hz）\\[1pt] LIPM / DCM / SRBD / Centroidal\\ 输出：期望接触力 $\mathbf{f}_i^\star$ $+$ 期望质心轨迹};
\node[ctrl, below=of plan] (ctrl) {\textbf{控制层：全身控制 WBC}（$\sim\!200$–$1000$\,Hz）\\[1pt] 任务优先级 QP / 逆动力学\\ 输出：关节力矩 $\boldsymbol{\tau}$ / 关节指令};
\node[motor, below=of ctrl] (motor) {\textbf{底层：电机驱动}（$\sim\!1$–$10$\,kHz）\\[1pt] 力矩 PD / 电流环};
\draw[ln] (plan) -- (ctrl) node[midway, right=1mm, font=\scriptsize]{期望力/运动};
\draw[ln] (ctrl) -- (motor) node[midway, right=1mm, font=\scriptsize]{力矩指令};
\end{tikzpicture}
\caption{腿足控制栈的三级金字塔。降阶模型只活在最上层；本章是这一层的“地图”。}
\label{fig:rm-pyramid}
\end{figure}

\begin{table}[H]\centering\small
\caption{降阶模型谱系总览——本章的主线锚（后续每节回扣此表）}
\label{tab:rm-landscape}
\begin{tabular}{lccl}
\toprule
模型 & 状态维度 & 核心假设 & 捕获的物理 \\
\midrule
LIPM & $2$–$4$ & 恒定 CoM 高度、点质量 & 线性平动 \\
DCM & $2$ & LIPM $+$ 发散分量分离 & 一阶不稳定性 \\
SLIP & $4$ & 弹簧腿、飞行相 & 垂直弹跳能量循环 \\
SRBD & $12$–$13$ & 整机单刚体、恒定 $\mathbf{I}_G$ & 平动 $+$ 转动 \\
LL-SRBM & $12$–$13$ & 单刚体 $+$ 腿质量聚合 & 平动 $+$ 转动（变 $\mathbf{I}_G$） \\
Centroidal & $12$ & 精确质心动量 & 完整 6D 动量 \\
\bottomrule
\end{tabular}
\end{table}

\begin{pitfall}[简化模型只用于规划层]
最常见的初学者错误，是试图\textbf{直接用降阶模型的输出去驱动电机}。降阶模型给出的是“期望质心力/接触力”，不是“关节力矩”——中间还隔着一层 WBC。LIPM 没有关节、SRBD 把整机当一块砖，它们\emph{物理上无法}回答“某个膝关节该出多大力矩”。把谱系当地图：它告诉你往哪走，但开车的是 WBC。
\end{pitfall}

\begin{exercise}
\begin{enumerate}
  \item Atlas 人形机器人约有 $28$ 个关节。仿照 \cref{eq:rm-decision-count} 估算其全身 MPC（$N=20$）的决策变量数量，并据 \cref{tab:rm-dim} 的趋势判断可行性。
  \item 为什么跑步的腾空期（飞行相）不能用 LIPM 建模？需要谱系里的哪个模型？（提示：飞行相破坏了哪条假设？）
  \item 用一句话解释：为什么“规划层 $+$ 控制层”的两级分工，比“单层全身 MPC”更可能满足实时性？
\end{enumerate}
\end{exercise}

% ════════════════════════════════════════════════════════════════
\section{LIPM：降阶谱系的起点}
\label{sec:rm-lipm}

\paragraph{物理设置。}
谱系的起点是最简单的双足模型：一个点质量 $m$ 代表机器人全部质量（集中在质心 $\mathrm{CoM}$），一条无质量刚性支撑腿连接 CoM 与地面支撑点 $p$，并\emph{假设 CoM 高度恒定} $z=h$（不考虑上下起伏）。这就是\textbf{线性倒立摆模型}（Linear Inverted Pendulum Mode, LIPM），由 Kajita 与 Tani 于 1991 年首次用于步行机器人\cite{kajita1991study}。

\paragraph{从牛顿定律到 LIPM 方程。}
质量只受重力 $m\mathbf{g}$ 与支撑腿约束力 $\mathbf{F}$。竖直方向因 $z=h$ 恒定（竖直加速度为零）给出 $F_z=mg$。支撑力沿腿方向（由 $p$ 指向 $m$），其水平分量在小角度（$h/l\approx 1$，$l$ 为腿长）下为 $F_x=mg\,(x-p)/h$。代入水平方向牛顿第二定律 $m\ddot{x}=F_x$：
\begin{equation}
  \boxed{\;\ddot{x}=\frac{g}{h}(x-p)=\omega_n^2\,(x-p)\;},\qquad \omega_n=\sqrt{g/h},
  \label{eq:rm-lipm}
\end{equation}
其中 $\omega_n$ 是 LIPM 的\textbf{自然频率}。因控制输入 $p$（CoP 位置）线性进入，\cref{eq:rm-lipm} 是关于 CoM 位置 $x$ 的\emph{线性} ODE——这正是 LIPM 名字里“Linear”的由来，也是它能引出后续所有解析结果的根。

\paragraph{本征不稳定性。}
齐次解（$p=0$）为 $x(t)=A\,e^{\omega_n t}+B\,e^{-\omega_n t}$，含两个模态：\emph{发散模态} $e^{\omega_n t}$（指数增长，对应摔倒）与\emph{收敛模态} $e^{-\omega_n t}$（指数衰减，自稳定）。任何微小扰动都会激发发散模态，故倒立摆\emph{本征不稳定}，必须主动控制。

\begin{insight}
LIPM 的不稳定性，就像在指尖上立一根棍子：你必须不断调整手指位置（CoP $p$）去抵消棍子（CoM $x$）的发散趋势，一旦停控，棍子立刻倒下。轮式机器人则像把球放在碗底——不控也会自然回到平衡点。这一根本差异，正是“为什么双足/腿足行走需要\emph{实时}控制”的物理答案。取 $h=0.8$\,m 得 $\omega_n=\sqrt{9.81/0.8}\approx 3.5$\,rad/s，则 $0.5$\,s 内扰动被 $e^{1.75}\approx 5.8$ 倍放大——不到一秒就会倒地。
\end{insight}

\paragraph{状态空间与离散化。}
取状态 $\mathbf{x}=(x,\dot{x})^\top$、控制 $u=p$：
\begin{equation}
  \dot{\mathbf{x}}=\begin{pmatrix}0&1\\\omega_n^2&0\end{pmatrix}\mathbf{x}
  +\begin{pmatrix}0\\-\omega_n^2\end{pmatrix}u,
  \label{eq:rm-lipm-ss}
\end{equation}
其系统矩阵特征值为 $\pm\omega_n$（一正一负），再次确认不稳定。MPC 需离散形式；零阶保持（采样 $\Delta t$）下矩阵指数给出闭式解
\begin{equation}
  \mathbf{A}_d=\begin{pmatrix}\cosh(\omega_n\Delta t)&\tfrac{1}{\omega_n}\sinh(\omega_n\Delta t)\\[2pt]
  \omega_n\sinh(\omega_n\Delta t)&\cosh(\omega_n\Delta t)\end{pmatrix},\qquad
  \mathbf{B}_d=\begin{pmatrix}1-\cosh(\omega_n\Delta t)\\-\omega_n\sinh(\omega_n\Delta t)\end{pmatrix}.
  \label{eq:rm-lipm-disc}
\end{equation}
（双曲函数取代振荡的三角函数，恰是“不稳定二阶系统”的指纹。）作为示意，离散矩阵在代码里仅两行：
\begin{codebox}[C++]{LIPM 零阶保持离散矩阵（示意）}
Ad << cosh(w*dt), sinh(w*dt)/w,  w*sinh(w*dt), cosh(w*dt);
Bd << 1.0 - cosh(w*dt),          -w*sinh(w*dt);
\end{codebox}

\paragraph{历史与局限。}
LIPM 的发展脉络是整条谱系的缩影：1991 年 Kajita–Tani 首次将其用于不平地形步行\cite{kajita1991study}；2001 年扩展为 3D-LIPM（增加 $y$ 方向，因 \cref{eq:rm-lipm} 在 $x,y$ 上解耦，故仍是两个独立的二阶 ODE）\cite{kajita2001lipm}；2003 年 Kajita 用预览控制（\cref{sec:rm-preview}）在 HRP-2 上实现动态行走。其代价是一组刚性假设，\cref{tab:rm-lipm-limit} 列出局限与对应的“升档”出口。

\begin{table}[H]\centering\small
\caption{LIPM 的局限与谱系中的升档出口}
\label{tab:rm-lipm-limit}
\begin{tabular}{lll}
\toprule
局限 & 后果 & 升档到 \\
\midrule
恒定 CoM 高度 & 不能跳跃／奔跑（有飞行相） & SLIP（\cref{sec:rm-slip}） \\
忽略角动量 & 不能甩臂／翻转 & SRBD/Centroidal（\cref{sec:rm-srbd-centroidal}） \\
点质量 & 不能预测关节力矩 & WBC（\cref{ch:control}） \\
平面运动 & $x,y$ 解耦、无耦合姿态 & 3D-LIPM（简单扩展） \\
\bottomrule
\end{tabular}
\end{table}

\begin{insight}
把 LIPM 看作谱系的\textbf{公共地基}：它提供了贯穿全章的 $\omega_n=\sqrt{g/h}$，是 ZMP（\cref{sec:rm-zmp}）、DCM（\cref{sec:rm-dcm}）、捕获点乃至 Raibert 落脚律（\cref{sec:rm-slip}）共同回扣的“二阶不稳定 ODE 范式”。理解了 \cref{eq:rm-lipm} 的两个模态，后面所有降阶模型都只是在“保留哪些模态、补回哪些被丢的物理”上做文章。
\end{insight}

\begin{note}
LIPM 在双足语境下的\emph{工程化}——DCM 步态反向递推、变高度 DCM、ALIP 横向落脚闭式律——已在 \cref{ch:legged-biped} 详尽展开。本章把 LIPM 定位为“谱系起点的范式推导”，DCM 的步态规划与反馈律见双足章，不在此重复。
\end{note}

\begin{exercise}
\begin{enumerate}
  \item 对 $y$ 方向重复 \cref{eq:rm-lipm} 的推导，验证 3D-LIPM 中 $x,y$ 完全解耦（两个独立二阶 ODE）。
  \item 设 $h=0.8$\,m，$\Delta t=10$\,ms。给定初值 $x_0=0,\dot{x}_0=0.1$\,m/s 且 CoP 固定在原点，用 \cref{eq:rm-lipm-disc} 解析估算 $2$\,s 后 $x$ 的量级（提示：$x(t)=(\dot{x}_0/\omega_n)\sinh(\omega_n t)$），并解释为什么是指数发散。
  \item 为什么四足机器人（如 Go2）尽管有 $4$ 条腿，仍能用 LIPM？（提示：考虑“等效支撑点”。）
\end{enumerate}
\end{exercise}

% ════════════════════════════════════════════════════════════════
\section{CoM／CoP／ZMP：三个关键点与稳定判据}
\label{sec:rm-zmp}

承上：LIPM 用单个支撑点 $p$ 概括了地面对机器人的作用。要把这个 $p$ 与真实的地面反力联系起来，需要精确区分三个易混的点——质心、压力中心、零力矩点。

\begin{definition}[质心 CoM]
多刚体系统的质心是各连杆质量加权的位置平均：
$\mathbf{r}_{\mathrm{CoM}}=\big(\sum_i m_i\mathbf{r}_i\big)\big/\big(\sum_i m_i\big)$。
\end{definition}

\begin{definition}[压力中心 CoP]
压力中心是地面反力（GRF）合力的等效作用点——在该点，地面对机器人的总力矩水平分量为零。平地点接触下
\begin{equation}
  \mathbf{p}_{\mathrm{CoP}}=\frac{\sum_i \mathbf{p}_i\, f_{z,i}}{\sum_i f_{z,i}},
  \label{eq:rm-cop}
\end{equation}
其中 $\mathbf{p}_i$ 是第 $i$ 个接触点、$f_{z,i}$ 是该点法向力。直觉：站在板上时，CoP 就是板对你施力的“加权平均位置”。
\end{definition}

\begin{definition}[零力矩点 ZMP]
零力矩点是地面上使地面合力矩水平分量为零的点。ZMP 概念由 Vukobratović 与 Juričić 于 \textbf{1968 年}（莫斯科 ICTAM 全苏理论与应用力学大会）首次提出，“零力矩点”这一术语在 \textbf{1970–1972 年间}的后续工作中确立\cite{vukobratovic2004zero}。
\end{definition}

\begin{theorem}[平地 ZMP$=$CoP]
\label{thm:rm-zmp-cop}
当所有接触点共面（平地行走）时，$\mathbf{p}_{\mathrm{ZMP}}=\mathbf{p}_{\mathrm{CoP}}$。
\end{theorem}
\begin{proof}[证明骨架]
平地上所有接触力均为竖直法向力，水平力矩只由竖直力的水平力臂产生。CoP 的定义 \cref{eq:rm-cop} 恰使这些水平力矩之和为零——这与 ZMP“水平力矩为零”的定义重合，故二者同点。
\end{proof}

\begin{theorem}[ZMP 稳定判据\cite{vukobratovic2004zero}]
\label{thm:rm-zmp-stable}
若 ZMP 位于支撑多边形内部，机器人不会绕支撑边界翻倒。三态判别：
\begin{itemize}
  \item ZMP 在多边形\emph{内部} $\Rightarrow$ 存在合法（非负法向力）的接触力分布可平衡机器人 $\Rightarrow$ 不翻倒；
  \item ZMP 在多边形\emph{边界} $\Rightarrow$ 某些脚的法向力趋零（即将抬起）$\Rightarrow$ 临界；
  \item ZMP 在多边形\emph{外部} $\Rightarrow$ 不存在合法接触力分布 $\Rightarrow$ 必然翻倒。
\end{itemize}
\end{theorem}

\begin{figure}[H]
\centering
\begin{tikzpicture}[>=Stealth, font=\small]
  % 支撑多边形（双脚）
  \foreach \x in {0,3.4}{
    \draw[fill=main!8, draw=main!60!black, rounded corners=1pt] (\x,0) rectangle ++(1.1,2.4);
  }
  \draw[densely dashed, draw=second!70!black, thick]
    (0,0) -- (4.5,0) -- (4.5,2.4) -- (0,2.4) -- cycle
    node[pos=0.62, above=8mm, fill=white, inner sep=1pt]{支撑多边形（凸包）};
  % 三态 ZMP
  \fill[third!80!black] (2.25,1.2) circle (2.2pt) node[below=1.5mm]{内部（稳定）};
  \fill[orange!90!black] (4.5,0.8) circle (2.2pt) node[right=1mm]{边界（临界）};
  \fill[red!85!black] (5.7,1.6) circle (2.2pt) node[right=1mm]{外部（翻倒）};
\end{tikzpicture}
\caption{ZMP 三态：相对支撑多边形（两脚凸包）的内部／边界／外部。}
\label{fig:rm-zmp}
\end{figure}

\paragraph{含角动量项的完整 ZMP 公式。}
从 CoM 处的牛顿–欧拉方程出发（重力 $+$ 惯性力的合力矩在 ZMP 处水平分量为零），可解出 ZMP 位置。记 $\mathbf{L}=(L_x,L_y,L_z)$ 为绕 CoM 的角动量，$m$ 为总质量，$z=z_{\mathrm{CoM}}$：
\begin{equation}
  p_{\mathrm{ZMP},x}=x_{\mathrm{CoM}}-\frac{z\,\ddot{x}_{\mathrm{CoM}}+\dot{L}_y/m}{g+\ddot{z}_{\mathrm{CoM}}},
  \qquad
  p_{\mathrm{ZMP},y}=y_{\mathrm{CoM}}-\frac{z\,\ddot{y}_{\mathrm{CoM}}-\dot{L}_x/m}{g+\ddot{z}_{\mathrm{CoM}}}.
  \label{eq:rm-zmp-full}
\end{equation}
两个分量的角动量项符号相反（$x$ 取 $+\dot{L}_y$、$y$ 取 $-\dot{L}_x$）：以地面原点对重力－惯性合扳手取矩 $\boldsymbol{\tau}^{gi}_O=\overrightarrow{OG}\times m\mathbf{g}-\dot{\mathbf{L}}_G$，再由 $\overrightarrow{OZ}=(\mathbf{n}\times\boldsymbol{\tau}^{gi}_O)/(\mathbf{f}^{gi}\!\cdot\!\mathbf{n})$、$\mathbf{n}=\mathbf{e}_z$ 展开即得（其中 $\mathbf{n}\times$ 把 $\dot{L}_y$ 送入 $x$ 分量、$\dot{L}_x$ 送入 $y$ 分量，分母 $\mathbf{f}^{gi}\!\cdot\!\mathbf{n}=-m(g+\ddot{z})$ 的负号再翻转一次符号）\cite{vukobratovic2004zero}。\footnote{此符号约定与 Tedrake 的 \emph{Underactuated Robotics}（其平面式写作 $(m\ddot z+mg)(x_{\mathrm{cp}}-x)=(z_{\mathrm{cp}}-z)m\ddot x-\dot{L}$）一致：当角动量增率 $\dot{\mathbf{L}}_G=0$ 时退化为下式 \cref{eq:rm-zmp-lipm}，正负不影响该退化极限。}注意量纲：$\dot{L}/m$ 单位为 $\mathrm{m^2/s^2}$，与 $z\ddot{x}$ 一致。

\begin{insight}
\cref{eq:rm-zmp-full} 在 LIPM 假设（$z=h$ 恒定、$\ddot{z}=0$、$\dot{L}=0$）下退化为
\begin{equation}
  p_{\mathrm{ZMP}}=x_{\mathrm{CoM}}-\frac{h}{g}\,\ddot{x}_{\mathrm{CoM}},
  \label{eq:rm-zmp-lipm}
\end{equation}
这恰是 LIPM 方程 \cref{eq:rm-lipm} 的另一种写法（代入 $\ddot{x}=\omega_n^2(x-p)$ 即得 $p_{\mathrm{ZMP}}=p_{\mathrm{CoP}}=p$）。\textbf{LIPM 的 CoP 控制，本质上就是 ZMP 控制}——这把上一节的支撑点 $p$ 与本节的地面反力焊在了一起。\cref{eq:rm-zmp-full} 还顺带告诉我们：角动量的变化率 $\dot{L}$（甩臂、摆腿）能把 ZMP 拉离 CoP，这正是后面 SRBD/Centroidal 要补回的转动物理。
\end{insight}

\begin{pitfall}[ZMP 判据是必要非充分]
ZMP 落在支撑多边形内，只保证“不绕支撑边界\emph{翻倒}”，\textbf{不保证}“不\emph{摔倒}”。滑动（摩擦不足）、关节力矩饱和、足端打滑等其它失效模式仍可能发生。换言之，ZMP $\in$ 支撑多边形是稳定行走的\emph{必要}条件而非\emph{充分}条件——它管住了一种失效，而非全部。
\end{pitfall}

\begin{exercise}
\begin{enumerate}
  \item 从 CoM 处的牛顿–欧拉力矩平衡出发，完整推出 \cref{eq:rm-zmp-full}，并核对两个角动量项的符号。
  \item 验证 \cref{eq:rm-zmp-full} 在 LIPM 假设下退化为 \cref{eq:rm-zmp-lipm}，并代入 \cref{eq:rm-lipm} 证明此时 $p_{\mathrm{ZMP}}=p_{\mathrm{CoP}}$。
  \item 举一个“ZMP 在支撑多边形内但机器人仍然摔倒”的具体场景，并指出它违反了哪条隐含假设。
\end{enumerate}
\end{exercise}

% ════════════════════════════════════════════════════════════════
\section{ZMP 预览控制（Kajita 2003）}
\label{sec:rm-preview}

承上：\cref{thm:rm-zmp-stable} 告诉我们“ZMP 该落在哪”，但没说“CoM 该怎么走才能让 ZMP 落在那”。这是一个\emph{逆问题}：给定步态规划器生成的期望 ZMP 轨迹 $p_{\mathrm{ref}}(t)$，求 CoM 轨迹 $x(t)$ 使实际 ZMP 跟踪参考。Kajita 2003 的预览控制是第一个实用答案\cite{kajita2001lipm}，也是全书唯一系统讲“预览（preview）”的地方。\footnote{此处的预览控制原始文献是 Kajita, Kanehiro, Kaneko, Fujiwara, Harada, Yokoi, Hirukawa, ``Biped Walking Pattern Generation by using Preview Control of Zero-Moment Point'', ICRA 2003, Taipei, pp.\ 1620--1626；它与 \cite{kajita2001lipm}（2001 IROS 的 3D-LIPM 建模）是两篇不同论文，前者方为预览控制（Cart-Table $+$ ZMP 跟踪伺服）之源。}

\paragraph{Cart-Table 模型。}
Kajita 用一个等价直觉——“推车–桌子”模型：质量 $m$ 的小车在桌上运动，桌子 ZMP 即桌面对小车的等效施力点。其状态方程以 CoM 的\emph{加加速度}（jerk）$u=\dddot{x}$ 为输入，构成三阶积分链：
\begin{equation}
  \frac{d}{dt}\begin{pmatrix}x\\\dot{x}\\\ddot{x}\end{pmatrix}
  =\begin{pmatrix}0&1&0\\0&0&1\\0&0&0\end{pmatrix}\begin{pmatrix}x\\\dot{x}\\\ddot{x}\end{pmatrix}
  +\begin{pmatrix}0\\0\\1\end{pmatrix}u,\qquad
  p=\begin{pmatrix}1&0&-h/g\end{pmatrix}\begin{pmatrix}x\\\dot{x}\\\ddot{x}\end{pmatrix}.
  \label{eq:rm-carttable}
\end{equation}
输出方程正是 ZMP 退化式 \cref{eq:rm-zmp-lipm}。与 LIPM \cref{eq:rm-lipm-ss} 相比，这里把 jerk 当输入、把 ZMP 当输出，从而能对 ZMP 直接施加跟踪与平滑代价。

\paragraph{增广系统与预览控制律。}
离散化（采样 $\Delta t$）后定义跟踪误差 $e_k=p_k-p_{\mathrm{ref},k}$，并把误差积分并入状态，构造增广状态 $\tilde{\mathbf{x}}_k=(e_k,\,x_k,\,\dot{x}_k,\,\ddot{x}_k)^\top$，增广动力学为 $\tilde{\mathbf{x}}_{k+1}=\tilde{\mathbf{A}}\tilde{\mathbf{x}}_k+\tilde{\mathbf{B}}u_k+\tilde{\mathbf{F}}\,p_{\mathrm{ref},k+1}$。在二次代价
\begin{equation}
  J=\sum_{k=0}^{\infty}\big[Q_e\,e_k^2+R\,u_k^2\big]
  \label{eq:rm-preview-cost}
\end{equation}
下求解（这是一个标准 LQR 问题，Riccati 理论见 \cref{ch:lqr-lqg-riccati}），最优控制律含三项：
\begin{equation}
  u_k=-K_I\sum_{i=0}^{k}e_i-\mathbf{K}_x\,\mathbf{x}_k-\sum_{j=1}^{N_p}G_j\,p_{\mathrm{ref},k+j},
  \label{eq:rm-preview-law}
\end{equation}
分别是误差积分项（$K_I$）、状态反馈项（$\mathbf{K}_x$）、与\textbf{预览项}（预览增益 $G_j$ 加权未来 $N_p$ 步的参考 ZMP）。预览增益由离散 Riccati 解 $\mathbf{P}$ 给出，
\begin{equation}
  G_j=(R+\tilde{\mathbf{B}}^\top\mathbf{P}\tilde{\mathbf{B}})^{-1}\tilde{\mathbf{B}}^\top
  (\tilde{\mathbf{A}}-\tilde{\mathbf{B}}\mathbf{K})^{j-1}\mathbf{P}\tilde{\mathbf{F}}.
  \label{eq:rm-preview-gain}
\end{equation}

\begin{insight}
\cref{eq:rm-preview-gain} 中 $(\tilde{\mathbf{A}}-\tilde{\mathbf{B}}\mathbf{K})^{j-1}$ 是\emph{闭环}状态转移矩阵的幂。因闭环稳定（谱半径 $<1$），$G_j$ 随预览步 $j$ \textbf{指数衰减}——越远的未来对当前决策影响越小。实践中取预览窗口 $N_p\approx 1.5$\,s 即足够（更远的 $G_j\approx 0$）。这解释了“为什么预知未来如此有用、又为什么不必预知太远”：脚步是提前规划好的，把未来 ZMP 喂给控制器能让 CoM \emph{提前}启动（避免临到换脚才手忙脚乱），但地面摩擦/反馈足以吸收一秒半以外的偏差。
\end{insight}

\paragraph{历史意义。}
1991 年之前，双足机器人只能静态站立或极慢行走；2003 年 Kajita 用 LIPM $+$ LQR $+$ preview 在 HRP-2 上首次实现动态行走\cite{kajita2001lipm}，证明了“一个二维线性模型 $+$ 一个线性二次调节器 $+$ 一个预览窗口”足以支撑实用的动态步行。这是降阶模型威力的第一个里程碑式注脚。

\begin{note}
\cref{eq:rm-preview-cost}--\cref{eq:rm-preview-gain} 的 LQR $+$ 预览理论与离散 Riccati 求解属于线性最优控制的一般结论，其代数与收敛性见 \cref{ch:lqr-lqg-riccati}。本节聚焦“ZMP 预览”的腿足应用与 Cart-Table 建模；通用预览控制不在此展开。
\end{note}

\begin{exercise}
\begin{enumerate}
  \item 由 \cref{eq:rm-carttable} 离散化并加入误差积分，推出增广矩阵 $\tilde{\mathbf{A}},\tilde{\mathbf{B}},\tilde{\mathbf{F}}$ 的表达式。
  \item 解释为什么 $R\ll Q_e$（强烈惩罚 ZMP 跟踪误差、轻惩罚 jerk）会给出更“贴”参考的 CoM 轨迹，代价是什么？
  \item 把预览窗口从 $1.5$\,s 缩到 $0$（纯反馈、无预览），定性说明 ZMP 跟踪误差会怎样变化，并联系 \cref{eq:rm-preview-gain} 的指数衰减解释。
\end{enumerate}
\end{exercise}

% ════════════════════════════════════════════════════════════════
\section{DCM 与捕获点：一阶降阶的本质}
\label{sec:rm-dcm}

承上：预览控制把 LIPM 的二阶系统当整体处理。本节换一个视角——把这个二阶系统\emph{分解}成稳定与不稳定两个一阶子系统，只盯住不稳定那一半。这正是发散分量（DCM）的思想，也是 M 谱系视角最有价值的洞察之一。

\paragraph{从 LIPM 到 DCM。}
回顾 LIPM 解 $x(t)=A\,e^{\omega_n t}+B\,e^{-\omega_n t}$。定义\textbf{发散分量}（Divergent Component of Motion, DCM）
\begin{equation}
  \xi=x+\frac{\dot{x}}{\omega_n}.
  \label{eq:rm-dcm-def}
\end{equation}
对 \cref{eq:rm-dcm-def} 求导并代入 LIPM 方程 \cref{eq:rm-lipm}：
\begin{equation}
  \dot{\xi}=\dot{x}+\frac{\ddot{x}}{\omega_n}=\dot{x}+\omega_n(x-p)
  =\omega_n\Big(x+\frac{\dot{x}}{\omega_n}\Big)-\omega_n p
  =\boxed{\;\omega_n(\xi-p)\;}.
  \label{eq:rm-dcm-dyn}
\end{equation}
\textbf{这是一阶线性 ODE}——LIPM 的二阶系统被降成了一阶。$\xi$ 恰是齐次解中\emph{发散模态的系数}，完全编码了系统的不稳定成分。

\paragraph{物理意义与捕获点。}
DCM $\xi$ 代表“若此刻停控（$p=0$），CoM 最终会发散到哪里”。由 \cref{eq:rm-dcm-dyn}：$\xi>p\Rightarrow\dot{\xi}>0$、$\xi<p\Rightarrow\dot{\xi}<0$，故\textbf{DCM 总是远离 CoP}——这是倒立摆不稳定性的一阶表达。\emph{捕获点}（Capture Point，Pratt et al.\ 2006\cite{pratt2006capture}）即当前 DCM 值 $\xi$：若机器人立即把脚踩到 $\xi$ 处（令 $p_{\mathrm{new}}=\xi$），则 $\dot{\xi}=\omega_n(\xi-\xi)=0$，DCM 停止发散，CoM 随后以 $e^{-\omega_n t}$ 收敛到 $\xi$，机器人停下不摔倒。据此可分三态：
\begin{itemize}
  \item 捕获点在可达范围内 $\Rightarrow$ 一步即可恢复平衡；
  \item 捕获点超出单步可达 $\Rightarrow$ 须多步恢复（进入“步态”）；
  \item 捕获点远超所有可能步态 $\Rightarrow$ 不可恢复，必然摔倒。
\end{itemize}
这一“可捕获性”分析框架由 Koolen et al.\ 2012 系统化\cite{koolen2012capturability}。

\begin{insight}
DCM 揭示了一个普适的控制原理：\textbf{对不稳定系统，控制资源应集中在不稳定模态上}。LIPM 有一个稳定模态（$e^{-\omega_n t}$，自动衰减、不必管）和一个不稳定模态（$e^{+\omega_n t}$，必须控制）。DCM 精确隔离了不稳定模态，使控制器“对症下药”——你只需把发散分量 $\xi$ 镇定，收敛分量会自行消失。这就是\emph{反馈镇定}的本质：任何线性系统都可分解为稳定与不稳定子空间，只需对不稳定子空间施加反馈。从谱系角度看，\textbf{DCM 就是 LIPM 的“一阶降阶”}——它丢掉了无须控制的收敛模态，把规划维度从二维（$x,\dot{x}$）压到一维（$\xi$）。
\end{insight}

\paragraph{收敛分量。}
与 DCM 对偶地，定义\textbf{收敛分量}（convergent component of motion, CCM）
\begin{equation}
  \zeta=x-\frac{\dot{x}}{\omega_n},\qquad \dot{\zeta}=-\omega_n(\zeta-p),
  \label{eq:rm-ccm}
\end{equation}
它自然衰减，即使不控制也不发散——这正是“控制器只须关注 DCM”的另一面。\footnote{部分文献会引入“虚拟排斥点”（Virtual Repellent Point, VRP，Englsberger et al.\ 2015\cite{englsberger2015dcm}）。须注意：VRP 属于\emph{排斥}结构（编码作用在机器人上的总力，DCM 被 VRP 排斥、CoM 被 DCM 吸引），\textbf{并非}此处的收敛分量 $\zeta$；二者不可混为一谈。}

\begin{pitfall}[变高度场景不要套用恒高 DCM]
DCM 公式 \cref{eq:rm-dcm-def} 中的 $\omega_n=\sqrt{g/h}$ 依赖恒定 CoM 高度 $h$。若机器人在不平地面行走（$h$ 变化），应改用\textbf{变高度 DCM}（time-varying DCM，自然频率随高度变化），或切换到\textbf{变高度倒立摆}（Variable-Height Inverted Pendulum, VHIP，Caron 2019/2020\cite{caron2019capturability}）这一标准模型。直接套用恒高 DCM 会在台阶、斜坡上系统性失准。
\end{pitfall}

\begin{note}
DCM 的\emph{工程化}——步态反向递推（从末步倒推每步起始 DCM）、DCM 反馈控制律、变高度 DCM 与 ALIP 横向落脚闭式律——已在 \cref{ch:legged-biped}（含 DCM 与捕获点的完整展开）做透。本节把 DCM 压成“谱系定位 $+$ 一阶降阶洞察”，不重铺其步态规划算法与反馈律。
\end{note}

\begin{exercise}
\begin{enumerate}
  \item 由 DCM 定义 \cref{eq:rm-dcm-def} 直接整理出 $x=\xi-\dot{x}/\omega_n$（等价地 $\dot{x}=\omega_n(\xi-x)$），并解释其物理意义（“CoM 速度恒指向 DCM，故 CoM 总在追随 DCM”）。\footnote{注意不可把此式误写成 $x=\xi-\dot{\xi}/\omega_n$：由 \cref{eq:rm-dcm-dyn} 有 $\dot{\xi}/\omega_n=\xi-p$，故 $\xi-\dot{\xi}/\omega_n=p$（恰是 CoP），而非 CoM 位置 $x$。}
  \item 求解一阶 ODE \cref{eq:rm-dcm-dyn} 得 $\xi(t)=p+(\xi_0-p)e^{\omega_n t}$；据此说明步态周期 $T_s$ 增大时步内发散量 $e^{\omega_n T_s}$ 如何变化，对步距规划意味着什么。
  \item 给出收敛分量 \cref{eq:rm-ccm} 的推导，并解释为什么“控制器只需关注 DCM”。
\end{enumerate}
\end{exercise}

% ════════════════════════════════════════════════════════════════
\section{SLIP：引入弹簧与飞行相}
\label{sec:rm-slip}

承上：LIPM 与 DCM 都建立在“CoM 高度恒定”之上，因而无法描述奔跑与跳跃。本节引入谱系中第一个\emph{非线性}模型——弹簧负载倒立摆（Spring-Loaded Inverted Pendulum, SLIP），它用弹簧腿与飞行相补回了被 LIPM 丢掉的垂直能量循环。SLIP 是全书唯一系统讲“弹簧倒立摆 $+$ Raibert”的地方。

\paragraph{为什么 LIPM 无法建模奔跑。}
奔跑与步行的本质区别是存在\textbf{飞行相}（所有脚离地）。飞行相中没有接触力，唯一的力是重力，CoM 做抛体运动 $\Rightarrow$ 高度变化 $\Rightarrow$ LIPM 的 $z=\mathrm{const}$ 假设破裂。SLIP 与 LIPM 的关系，恰如弹性力学与刚体力学：后者假设不变形（CoM 高度不变），前者引入弹性（弹簧腿允许 CoM 上下起伏），适用范围因此大幅扩展。

\paragraph{SLIP 物理模型与两相方程。}
点质量 $m$ 由一条弹簧腿（刚度 $k$、自然长 $L_0$）支撑，两相交替（\cref{fig:rm-slip}）：
\begin{itemize}
  \item \textbf{支撑相}（stance）：弹簧压缩–回弹，$m\ddot{\mathbf{r}}=-k(r-L_0)\hat{\mathbf{r}}+m\mathbf{g}$；
  \item \textbf{飞行相}（flight）：自由抛体，$m\ddot{\mathbf{r}}=m\mathbf{g}$。
\end{itemize}
以足端为原点、用极坐标 $(r,\theta)$（$\theta$ 自竖直方向量起）描述支撑相，径向加速度 $a_r=\ddot{r}-r\dot{\theta}^2$、切向加速度 $a_\theta=r\ddot{\theta}+2\dot{r}\dot{\theta}$，由径向（弹簧力 $+$ 重力径向分量）与切向（重力切向分量）力平衡得
\begin{equation}
  m(\ddot{r}-r\dot{\theta}^2)=-k(r-L_0)-mg\cos\theta,\qquad
  m(r\ddot{\theta}+2\dot{r}\dot{\theta})=mg\sin\theta.
  \label{eq:rm-slip}
\end{equation}
非线性来自离心项 $r\dot{\theta}^2$、科氏项 $2\dot{r}\dot{\theta}$ 与 $\cos\theta/\sin\theta$ 的耦合——\cref{eq:rm-slip} \textbf{没有闭式解}，只能数值积分（实践中用 RK4：在每相内用 $\mathbf{k}_1$–$\mathbf{k}_4$ 四次斜率估计推进一步，并在 $r=L_0$ 处检测相切换）。

\begin{figure}[H]
\centering
\begin{tikzpicture}[>=Stealth, font=\small]
  \draw[thick] (-0.3,0) -- (8.5,0); % ground
  % flight 1 (apex, projectile)
  \fill[main!80!black] (0.6,2.6) circle (3pt) node[above]{$m$};
  \draw[->, second!70!black, thick] (0.6,2.6) .. controls (1.2,1.6) and (1.6,0.9) .. (2.0,0.55);
  \node[font=\scriptsize] at (0.9,1.5){飞行相（抛体）};
  % stance (spring)
  \fill[main!80!black] (3.7,1.9) circle (3pt) node[above]{$m$};
  \draw[decorate, decoration={coil, segment length=2.4pt, amplitude=2.4pt}, third!60!black, thick]
    (3.2,0) -- (3.7,1.9);
  \fill[black] (3.2,0) circle (1.6pt);
  \node[font=\scriptsize] at (4.7,1.0){支撑相（弹簧 $k,L_0$）};
  % flight 2
  \fill[main!80!black] (7.0,2.6) circle (3pt) node[above]{$m$};
  \draw[->, second!70!black, thick] (5.6,0.55) .. controls (6.2,1.1) and (6.6,1.9) .. (7.0,2.6);
\end{tikzpicture}
\caption{SLIP 的飞行–支撑–飞行循环：飞行相做抛体运动（CoM 高度变化），支撑相由弹簧腿压缩回弹完成垂直能量交换。}
\label{fig:rm-slip}
\end{figure}

\paragraph{Poincaré 映射与被动自稳定。}
既无闭式解，分析 SLIP 稳定性的标准工具是 \textbf{Poincaré 映射}：取截面为 apex（最高点，$\dot{z}=0$），状态 $(z_{\mathrm{apex}},\dot{x}_{\mathrm{apex}})$，映射 $P$ 为“从一个 apex 到下一个 apex”的状态转移；不动点 $\mathbf{x}^\star=P(\mathbf{x}^\star)$ 即周期跳跃/奔跑。稳定性由 Jacobian $\mathbf{J}=\partial P/\partial\mathbf{x}|_{\mathbf{x}^\star}$ 的特征值判定：若所有 $|\lambda_i(\mathbf{J})|<1$ 则渐近稳定。SLIP 的一个惊人性质是\textbf{被动自稳定}——在特定参数范围内，\emph{固定落脚角}即可对小扰动自动恢复，无须主动反馈（Seyfarth/Geyer/Blickhan\cite{geyer2006compliant}），与被动行走机器人的自稳定异曲同工。

\paragraph{Raibert 三分解耦（1986）。}
Raibert 在没有 SLIP 理论分析的年代，凭工程直觉提出了惊人有效的解耦控制\cite{raibert1986legged}：
\begin{enumerate}
  \item \textbf{跳跃高度}：调节弹簧推力；
  \item \textbf{前进速度}：调节落脚点（Raibert 启发式）
    \begin{equation}
      p_{\mathrm{foot}}=\frac{v_x T_s}{2}+k_v(v_x-v_{\mathrm{ref}});
      \label{eq:rm-raibert}
    \end{equation}
  \item \textbf{身体姿态}：支撑相中用髋关节力矩保持身体水平。
\end{enumerate}
\cref{eq:rm-raibert} 中第一项 $v_x T_s/2$ 是\emph{中性点}（neutral point）——若落在此处，支撑相前半减速恰抵消后半加速，速度不变；第二项 $k_v(v_x-v_{\mathrm{ref}})$ 是反馈校正（太快则落脚偏前以多减速）。作为示意，这一启发式只是一行：
\begin{codebox}[C++]{Raibert 落脚启发式（示意）}
p_foot = 0.5 * vx * Ts + kv * (vx - v_ref);
\end{codebox}

\begin{insight}
\cref{eq:rm-raibert} 与捕获点（\cref{sec:rm-dcm}）有深层联系：$v_x T_s/2$ 是“假设匀减速时支撑相中点的预测位置”，本质上是\textbf{一阶近似的捕获点控制}；反馈项则对应 DCM 反馈。换言之，\emph{Raibert 启发式就是 SLIP 上的“简化 DCM 控制”}——这是谱系的一个漂亮连接点：看似各自独立的弹簧倒立摆与线性倒立摆，在落脚律层面殊途同归。Raibert 启发式至今仍是最常用的落脚规划公式，其鲁棒性来自只依赖速度 $v_x$ 与步态周期 $T_s$（不需精确质量/惯量/刚度）。
\end{insight}

\paragraph{3D-SLIP 与四足步态。}
把 2D-SLIP 扩展到 3D 后，四足的不同步态可看作不同的 SLIP 配置（\cref{tab:rm-slip-gait}）：trot 约等于两个对角 SLIP；bound 约等于纵向前后双 SLIP（需飞行相，故 LIPM 无法建模）；gallop 是非对称双 SLIP（前后落地存在相位差，比 bound 更高效）。MIT Cheetah 的 bounding/galloping 本质上就是 3D-SLIP。

\begin{table}[H]\centering\small
\caption{四足步态与 SLIP 等效（概念对照，作直觉）}
\label{tab:rm-slip-gait}
\begin{tabular}{llll}
\toprule
步态 & SLIP 等效 & 飞行相 & 典型平台 \\
\midrule
步行（walk） & LIPM（无弹簧） & 无 & HRP-2, Atlas \\
小跑（trot） & 双对角 SLIP & 无／短 & Go2, Spot \\
奔跑（bound） & 纵向双 SLIP & 有 & MIT Cheetah \\
跳跃（pronk） & 单 SLIP & 有 & ANYmal \\
\bottomrule
\end{tabular}
\end{table}

\begin{pitfall}[SLIP 的弹簧是等效刚度]
SLIP 的弹簧刚度 $k$ 不是物理弹簧，而是\textbf{等效刚度}，由腿的关节阻抗／力控策略产生。真实机器人通过 PD 力控或阻抗控制“模拟”弹簧行为（阻抗控制见 \cref{ch:control}）。把 $k$ 当成真实弹性元件去标定，会与实际腿控参数对不上。
\end{pitfall}

\begin{exercise}
\begin{enumerate}
  \item 由笛卡尔坐标下的弹簧力 $\mathbf{F}_s=-k(r-L_0)\hat{\mathbf{r}}$ 出发，完整推出极坐标支撑相方程 \cref{eq:rm-slip}，并指出三个非线性项各自的物理来源。
  \item 给定一条周期跳跃轨迹，写出数值计算 Poincaré 映射 Jacobian $\partial P/\partial\mathbf{x}$ 并据特征值判稳的步骤。
  \item 由“匀减速、支撑相中点”推出中性点 $v_x T_s/2$，并解释它为何近似等于一阶捕获点。
\end{enumerate}
\end{exercise}

% ════════════════════════════════════════════════════════════════
\section{SRBD 与 Centroidal：谱系的高动态两档}
\label{sec:rm-srbd-centroidal}

承上：LIPM/DCM/SLIP 都用点质量描述平动，丢掉了转动。要预测转弯、坡面俯仰横滚、四足 trot 的机身振荡，必须补回转动动力学。谱系在此进入两个高动态档：单刚体模型（SRBD）与质心动力学模型（Centroidal）。

\subsection{SRBD：把整机当一块砖}
\label{subsec:rm-srbd}

\paragraph{单刚体假设与两方程。}
SRBD 把整个机器人视为一个位于质心 $\mathbf{r}_G$、总质量 $m$、惯量 $\mathbf{I}_G$ 的刚体，\emph{忽略所有腿部质量与惯性}。其平动（牛顿）与转动（欧拉）方程为
\begin{align}
  m\ddot{\mathbf{r}}_G&=\sum_{i=1}^{n_c}\mathbf{f}_i+m\mathbf{g},
  \label{eq:rm-srbd-trans}\\
  \mathbf{I}_G\dot{\boldsymbol{\omega}}+\boldsymbol{\omega}\times(\mathbf{I}_G\boldsymbol{\omega})
  &=\sum_{i=1}^{n_c}(\mathbf{p}_i-\mathbf{r}_G)\times\mathbf{f}_i.
  \label{eq:rm-srbd-rot}
\end{align}
单刚体假设的物理合理性在于：四足机器人腿部质量通常仅占总质量 $10$–$15\%$，且电机常被置于近髋处以降低腿部转动惯量。平动方程 \cref{eq:rm-srbd-trans} 实际\emph{精确}（牛顿第二定律对质心恒成立，与是否刚体无关）；近似只在“$m$ 取常数、用简化 CoM、$\mathbf{I}_G$ 取常数”这几处。$\boldsymbol{\omega}\times(\mathbf{I}_G\boldsymbol{\omega})$ 是陀螺力矩，$\boldsymbol{\omega}$ 小时可忽略。

\begin{insight}
SRBD 的真正威力不在“简化”，而在\textbf{把决策变量从关节空间（$12$ 维力矩）转移到接触力空间（同样 $12$ 维力，但含义全然不同）}。接触力直接对应牛顿–欧拉方程的右端项——优化器能“看见”力如何影响质心运动，无须经过关节动力学这个中间层。这就是为什么把 \cref{eq:rm-srbd-trans}--\cref{eq:rm-srbd-rot} 线性化后能写成\emph{凸二次规划}、从而毫秒级求解的根本原因。
\end{insight}

\paragraph{谱系定位：SRBD 丢了什么。}
相对 LIPM，SRBD 补回了转动（\cref{eq:rm-srbd-rot}）；相对全身动力学，它丢掉了关节级动力学与\emph{惯量随构型的变化}（假设 $\mathbf{I}_G$ 恒定）。世界系惯量经右扰动姿态 $\mathbf{R}\in\SO(3)$ 变换为 $\mathbf{I}_G^W=\mathbf{R}\,\mathbf{I}_G^B\,\mathbf{R}^\top$，小角度线性化后成为姿态角的线性函数，这是“为何能凸”的关键一环。SRBD 与质心动量的关系是
\begin{equation}
  \mathbf{h}_G^{\mathrm{SRBD}}=\mathbf{A}_G^{[1:6,\,1:6]}(q)\,\dot{q}_{\mathrm{base}}
  \approx\begin{bmatrix}m\dot{\mathbf{r}}_G\\\mathbf{I}_G\boldsymbol{\omega}\end{bmatrix}
  \quad(\text{线动量在前}),
  \label{eq:rm-srbd-cmm}
\end{equation}
即只取质心动量矩阵 $\mathbf{A}_G$ 对应浮动基座的 $6\times 6$ 块、忽略关节速度对动量的贡献。当腿摆幅度大（跳跃、奔跑），被忽略的 $\mathbf{A}_G^{[\,\cdot\,,\,7:n_v]}\dot{q}_{\mathrm{joints}}$ 不可忽略，SRBD 精度显著下降——这正是升档到 Centroidal 的动机。

\begin{note}
SRBD 的\emph{完整推导}（密度积分、世界系惯量线性化）、\emph{三大凸化近似}（丢陀螺项、小角度欧拉角率映射、摩擦金字塔）、\emph{$13$ 维状态 $+$ 重力增广}、连续 $\mathbf{A}_c/\mathbf{B}_c$ 显式结构、前向欧拉离散、$N$ 步\emph{堆叠 QP}、计算量表、warm-start、以及 Di Carlo 2018 的全部工程细节，均已在 \cref{ch:quad_mpc} 逐一覆盖\cite{dicarlo2018cheetah}。本节只给谱系定位与上述 CMM 关系，\textbf{不重复其推导、状态构造、QP 堆叠与代码}。
\end{note}

\subsection{Centroidal：谱系顶端的精确档}
\label{subsec:rm-centroidal}

\paragraph{动机与质心动量动力学。}
SRBD 假设 $\mathbf{I}_G$ 恒定，但机器人大幅改变姿态（如人形甩臂）时 $\mathbf{I}_G$ 显著变化。\textbf{质心动力学}（Centroidal Dynamics）不做此假设，而直接对 6D 质心动量 $\mathbf{h}_G$ 列写动力学。按本书线动量在前的约定：
\begin{equation}
  \dot{\mathbf{h}}_G=\begin{pmatrix}\dot{\mathbf{l}}_G\\\dot{\mathbf{k}}_G\end{pmatrix}
  =\begin{pmatrix}\sum_i\mathbf{f}_i+m\mathbf{g}\\[2pt]\sum_i(\mathbf{p}_i-\mathbf{r}_G)\times\mathbf{f}_i\end{pmatrix},
  \label{eq:rm-centroidal-dyn}
\end{equation}
其中 $\mathbf{l}_G=m\dot{\mathbf{r}}_G$ 是线动量、$\mathbf{k}_G$ 是绕 CoM 的角动量。它与 CMM 的桥是 $\mathbf{h}_G=\mathbf{A}_G(q)\dot{q}$（\cref{sec:rm-bridge}）。\cref{tab:rm-srbd-vs-cen} 给出 SRBD 与 Centroidal 的数学分界。

\begin{table}[H]\centering\small
\caption{SRBD 与 Centroidal 的数学区别}
\label{tab:rm-srbd-vs-cen}
\begin{tabular}{lll}
\toprule
 & SRBD & Centroidal \\
\midrule
角动量表达 & $\mathbf{k}_G=\mathbf{I}_G^{\mathrm{const}}\boldsymbol{\omega}$ & $\mathbf{k}_G=\mathbf{A}_G^{\mathrm{ang}}(q)\,\dot{q}$ \\
$\mathbf{I}_G$ 假设 & 恒定（或仅依姿态角线性化） & 无假设（精确，随 $q$ 非线性） \\
优化结构 & 凸 QP（线性化后） & NLP（双线性 $+$ 构型依赖） \\
实时性 & 高 & 中（需 SQP/DDP） \\
\bottomrule
\end{tabular}
\end{table}

\paragraph{为何 Centroidal MPC 是非线性规划。}
两处非线性使 \cref{eq:rm-centroidal-dyn} 无法直接写成 QP：其一，叉积 $(\mathbf{p}_i-\mathbf{r}_G)\times\mathbf{f}_i$ 是质心位置与力的\emph{双线性}耦合；其二，$\mathbf{A}_G(q)$ 是关节角度的\emph{非线性}函数。常用三类求解法（概念对照，求解器理论分别见 \cref{ch:nlopt} 与 \cref{ch:ddp-ilqr}）：\textbf{SQP/SLQ}（逐步线性化解 QP，超线性收敛，自然支持约束，代表 OCS2）；\textbf{DDP/FDDP}（利用二阶信息，二次收敛，对初值敏感，代表 Crocoddyl）；\textbf{交替优化}（固定位置求力、固定力求轨迹，线性收敛）。三者都依赖 warm-start 把迭代次数从 $10$–$20$ 降到 $2$–$5$。

\begin{insight}
\textbf{Orin 2013 的关键洞察}\cite{orin2013centroidal}：$\dot{\mathbf{h}}_G$ 只依赖\emph{接触力}与\emph{质心位置}，不直接依赖关节角度。这意味着 Centroidal MPC 的动力学约束是“低维的”（$12$ 维），尽管完整系统有 $30^+$ 维；关节角度通过 $\mathbf{A}_G(q)$ 影响 $\mathbf{h}_G$，但这个耦合可放进\emph{代价函数}而非\emph{硬约束}。正是这一结构，让 Centroidal 成为“最精确却仍可处理”的简化模型——谱系的顶端。
\end{insight}

\begin{note}
CMM/CCRBA 的 $O(N)$ 算法、复合刚体质心惯量（CCRBI）$\mathbf{I}_G=\mathbf{A}_G\mathbf{M}^{-1}\mathbf{A}_G^\top$ 等，归 \cref{ch:contact-mechanics}（接触力学章已认领其算法主体）\cite{orin2013centroidal}；OCS2 kino-centroidal、Crocoddyl FDDP 等质心 NMPC 的工程实现见 \cref{ch:quad_mpc}。本节只用 $\mathbf{h}_G=\mathbf{A}_G(q)\dot{q}$ 这一线性桥作概念接口。SRBD 与 Centroidal 在 Go2 trot 上的定量精度对照（CoM 误差与角动量误差随预测时域增长）为 Pinocchio 仿真经验量级、依实现而变\rebuilt{无原始出处的具体百分比，仅作趋势直觉}：定性结论稳健——腿质量占比小故 SRBD 的 CoM 误差小（毫米级），但角动量误差随时域显著增大（可达两位数百分比），因为腿摆引起的角动量变化被 SRBD 完全忽略。
\end{note}

\begin{exercise}
\begin{enumerate}
  \item 由 \cref{eq:rm-srbd-rot} 出发，写出陀螺项 $\boldsymbol{\omega}\times(\mathbf{I}_G\boldsymbol{\omega})$ 可忽略的两个条件，并各举一个会使其\emph{不可}忽略的运动。
  \item 解释 \cref{eq:rm-srbd-cmm} 中“只取 $\mathbf{A}_G$ 的浮动基座 $6\times 6$ 块”等价于什么物理假设；在何种步态下该假设失效。
  \item 指出 \cref{eq:rm-centroidal-dyn} 的两处非线性来源，并据 \cref{tab:rm-srbd-vs-cen} 说明为什么它一般只能用 SQP/DDP 而非 QP 求解。
\end{enumerate}
\end{exercise}

% ════════════════════════════════════════════════════════════════
\section{模型谱系与选型理论}
\label{sec:rm-selection}

前六节逐档走完了谱系。本节把它们横向铺开，回答工程上最实际的问题：\textbf{面对一个具体任务，该选哪个模型？} 这是本章的理论高峰，也是 P7 各章“一章一模型”所没有的总图。

\paragraph{选型表与决策树。}
\cref{tab:rm-select} 按“维度／计算／旋转／跳跃／三类机器人”给出六模型的横向对照，\cref{fig:rm-decision} 则把选型逻辑画成一棵二叉决策树。

\begin{table}[H]\centering\small
\caption{降阶模型选型表（六模型横向对照；$\checkmark$ 可，$\triangle$ 受限，$\times$ 不可）}
\label{tab:rm-select}
\begin{tabular}{lccccccc}
\toprule
模型 & 状态维度 & 计算 & 旋转 & 跳跃 & 四足 & 双足 & 人形 \\
\midrule
LIPM       & $2$–$4$  & 极低 & $\times$ & $\times$ & $\triangle$ & $\checkmark$ & $\triangle$ \\
DCM        & $2$      & 极低 & $\times$ & $\times$ & $\triangle$ & $\checkmark$ & $\triangle$ \\
SLIP       & $4$      & 中   & $\times$ & $\checkmark$ & $\checkmark$ & $\triangle$ & $\times$ \\
SRBD       & $12$–$13$& 中（凸 QP） & $\checkmark$ & $\triangle$ & $\checkmark$ & $\triangle$ & $\triangle$ \\
Centroidal & $12$     & 高（NLP） & $\checkmark$ & $\checkmark$ & $\checkmark$ & $\checkmark$ & $\checkmark$ \\
全身动力学  & $2n$     & 极高 & $\checkmark$ & $\checkmark$ & $\triangle$离线 & $\triangle$离线 & $\triangle$离线 \\
\bottomrule
\end{tabular}
\end{table}

\begin{figure}[H]
\centering
\begin{tikzpicture}[
  >=Stealth, font=\small, node distance=7mm and 12mm,
  q/.style={draw, rounded corners, align=center, font=\small, inner sep=4pt, fill=main!8, draw=main!60!black},
  a/.style={draw, rounded corners, align=center, font=\small, inner sep=4pt, fill=third!10, draw=third!60!black},
  ln/.style={->, thick, gray!55!black}]
\node[q] (rt) {需要实时 MPC？};
\node[a, below left=of rt] (off) {全身动力学\\（离线规划／训练）};
\node[q, below right=of rt] (q2) {只需平动\\（步行）？};
\node[a, below left=12mm of q2] (lip) {LIPM / DCM\\（最简最快）};
\node[q, below right=of q2] (q3) {惯量变化大？};
\node[a, below left=10mm of q3] (srb) {SRBD\\（凸 QP）};
\node[a, below right=of q3] (cen) {Centroidal\\（SQP/DDP）};
\draw[ln] (rt) -- node[above left, font=\scriptsize]{否} (off);
\draw[ln] (rt) -- node[above right, font=\scriptsize]{是} (q2);
\draw[ln] (q2) -- node[above left, font=\scriptsize]{是} (lip);
\draw[ln] (q2) -- node[above right, font=\scriptsize]{否} (q3);
\draw[ln] (q3) -- node[above left, font=\scriptsize]{否} (srb);
\draw[ln] (q3) -- node[above right, font=\scriptsize]{是} (cen);
\end{tikzpicture}
\caption{降阶模型选型决策树：从“是否实时”到“惯量是否变化大”的二叉判别。}
\label{fig:rm-decision}
\end{figure}

\paragraph{信息丢失分析——谱系的灵魂。}
选型的关键不是“哪个最精确”，而是\textbf{理解每个模型抛弃了什么物理}——抛弃的若与任务无关则可放心降档，若与任务核心相关则必须升档。\cref{tab:rm-infoloss} 是本章的灵魂表。

\begin{table}[H]\centering\small
\caption{各降阶模型的“信息丢失”及其后果}
\label{tab:rm-infoloss}
\begin{tabular}{lll}
\toprule
模型 & 丢失的信息 & 后果 \\
\midrule
LIPM       & 旋转、CoM 高度变化、腿动力学 & 无法预测跌倒方向、无法奔跑 \\
DCM        & 同 LIPM $+$ 连续时间信息（离散化） & 步间无法修正 \\
SLIP       & 旋转、多体耦合 & 无法控制身体姿态 \\
SRBD       & 关节级动力学、惯量变化 & 大幅运动时力分配不准 \\
Centroidal & 关节级约束（力矩／速度限制） & 可能生成物理不可执行轨迹 \\
\bottomrule
\end{tabular}
\end{table}

\paragraph{LL-SRBM：谱系的中间档。}
标准 SRBD 把全部质量集中在躯干，对腿质量占比 $30$–$40\%$ 的人形或大型四足误差显著。\textbf{腿质量聚合单刚体}（Lump-Leg SRBM, LL-SRBM）是 SRBD 与 Centroidal 之间的中间档：把每条腿质量“聚集”到其质心，再用平行轴定理并入整机惯量，
\begin{equation}
  \mathbf{I}_G^{\mathrm{LL}}=\mathbf{I}_{\mathrm{body}}+\sum_{i=1}^{n_{\mathrm{legs}}}m_{\mathrm{leg},i}\,d_i^2,
  \label{eq:rm-llsrbm}
\end{equation}
其中 $d_i$ 是第 $i$ 条腿聚合质心到躯干质心的距离。状态空间与 SRBD 相同，但 $\mathbf{I}_G^{\mathrm{LL}}$ 随构型变化——它\emph{仍是凸 QP}（每个 MPC 步开始时用当前关节角更新 $\mathbf{I}_G^{\mathrm{LL}}$ 即可），是“计算几乎不增、精度显著提升”的升级选项，填补了谱系中“SRBD 太糙、Centroidal 太重”的空档。

\paragraph{量化误差框架与升降档规则。}
选型的定量依据是模型预测误差——简化模型与全身动力学（Pinocchio 真值）算出的 CoM 加速度之相对差：
\begin{equation}
  e_{\mathrm{model}}=\frac{\big\|\ddot{\mathbf{r}}_G^{\mathrm{simpl}}-\ddot{\mathbf{r}}_G^{\mathrm{full}}\big\|}
  {\big\|\ddot{\mathbf{r}}_G^{\mathrm{full}}\big\|}.
  \label{eq:rm-emodel}
\end{equation}
经验上的工程判断规则\rebuilt{以下百分比为代表性量级、依平台/步态/求解器而变，非普适基准}：$e_{\mathrm{model}}<5\%$ 时 MPC 反馈足以补偿、无须升级；$5\%$–$15\%$ 时 WBC 可补偿大部分但性能略降；$>15\%$ 时应升级模型或增大反馈增益。何时\emph{从 SRBD 升到 Centroidal}：腿质量占比 $>20\%$、运动中角动量变化显著（跳跃/翻转）、时域内有步态切换、或需规划手臂运动。何时\emph{从 Centroidal 回退 SRBD}：算力不足、运动以低速行走为主（$e_{\mathrm{model}}<5\%$）、或快速原型阶段（SRBD $+$ QP 更易调试）。LIPM/SRBD/Centroidal 的\emph{定性排序}（精度递增、算力递增）是稳健的。

\paragraph{谱系的下游接口。}
\cref{tab:rm-downstream} 总结各模型与下游控制框架（MPC/WBC/RL）的配合方式——这是谱系“向下连接”到控制层的接口图。

\begin{table}[H]\centering\small
\caption{各降阶模型与下游控制框架的配合}
\label{tab:rm-downstream}
\begin{tabular}{llll}
\toprule
模型 & 作为 MPC 的动力学约束 & 作为 WBC 的参考 & 作为 RL 的奖励参考 \\
\midrule
LIPM       & ZMP 约束 $+$ 预览 & 期望 CoP 轨迹 & 少用 \\
DCM        & 步态规划（反向递推） & 期望 DCM 轨迹 & 步距参考 \\
SRBD       & 凸 QP 约束（Di Carlo） & 期望接触力 $+$ CoM 轨迹 & 力参考 \\
Centroidal & SQP/DDP 约束（OCS2） & 期望动量 $+$ 接触力 & 动量参考 \\
\bottomrule
\end{tabular}
\end{table}

\begin{exercise}
\begin{enumerate}
  \item 用 \cref{fig:rm-decision} 的决策树，为下列任务各选一个模型并说明理由：(a) HRP-2 平地慢走；(b) Go2 平地 trot；(c) ANYmal 原地跳 $0.3$\,m；(d) 人形机器人甩臂平衡。
  \item 解释为什么 LL-SRBM（\cref{eq:rm-llsrbm}）仍是凸 QP，而 Centroidal 不是。（提示：$\mathbf{I}_G^{\mathrm{LL}}$ 在一个 MPC 步内是常数还是变量？）
  \item 据 \cref{tab:rm-infoloss}，说明“Centroidal 也会生成物理不可执行轨迹”的原因，以及哪一层（WBC？）负责把它拉回可行域。
\end{enumerate}
\end{exercise}

% ════════════════════════════════════════════════════════════════
\section{浮动基座 $\leftrightarrow$ 质心动量：全身与简化的精确桥}
\label{sec:rm-bridge}

承上：选型表把六个模型并列，但它们与\emph{全身}动力学究竟是什么关系？本节给出把整张谱系收束起来的“桥”——质心动量矩阵（CMM）。它是一座精确的投影桥，让“全身 $\to$ 简化”不再是模糊的类比，而是一个线性映射。

\paragraph{全身动力学：维度灾难的源头（引用，不重证）。}
\cref{ch:spatial-dynamics} 已建立浮动基座机器人的全身动力学方程
\begin{equation}
  \mathbf{M}(q)\ddot{q}+\mathbf{h}(q,\dot{q})=\mathbf{S}^\top\boldsymbol{\tau}+\mathbf{J}_c^\top\boldsymbol{\lambda},
  \label{eq:rm-fullbody}
\end{equation}
其中 $\mathbf{M}(q)\in\mathbb{R}^{n_v\times n_v}$ 是广义质量矩阵、$\mathbf{h}$ 是科氏／离心／重力项、选择矩阵 $\mathbf{S}=[\mathbf{0}_{n\times 6}\ \ \mathbf{I}_{n}]$ 把关节力矩映入广义力（故 $\mathbf{S}^\top\boldsymbol{\tau}$ 前 $6$ 个分量为零——基座\emph{欠驱动}），$\mathbf{J}_c$ 是接触 Jacobian、$\boldsymbol{\lambda}$ 是接触力。\cref{eq:rm-fullbody} 正是 \cref{sec:rm-curse} 维度灾难的源头（Go2 $n_v=18$）。其完整推导、选择矩阵的块分解（基座行／关节行）与接触 Jacobian 的构造均见 \cref{ch:spatial-dynamics}，本节只把它当作“被降阶的对象”引用一次。

\paragraph{质心动量桥。}
全身状态向简化模型的精确投影，由质心动量矩阵 $\mathbf{A}_G(q)\in\mathbb{R}^{6\times n_v}$ 实现：
\begin{equation}
  \boxed{\;\mathbf{h}_G=\mathbf{A}_G(q)\,\dot{q}=\begin{bmatrix}\mathbf{l}_G\\\mathbf{k}_G\end{bmatrix}\;}
  \qquad(\mathbf{l}_G=m\dot{\mathbf{r}}_G,\ \text{线动量在前}).
  \label{eq:rm-cmm}
\end{equation}
$\mathbf{A}_G$ 的每一列表示“对应广义速度分量变化一单位时，质心动量变化多少”。其精确性在于：它把\emph{所有}连杆（含腿）对总动量的贡献都算进来，而非只看基座。与之配套的质心 Newton–Euler 方程恰是 \cref{eq:rm-centroidal-dyn}——\textbf{$\dot{\mathbf{h}}_G$ 只等于外部接触力与重力在 CoM 处的合扳手}，与内部关节如何运动无关。这正是 \cref{eq:rm-fullbody} 的“基座 $6$ 行”（无 $\boldsymbol{\tau}$、只受接触力驱动）在质心坐标下的等价表述。

\begin{insight}
CMM \cref{eq:rm-cmm} 是理解整张谱系的钥匙。它把每个降阶模型都解释成“对 $\mathbf{A}_G$ 的某种近似”（\cref{tab:rm-cmm-bridge}）：LIPM 假设角动量恒为零（$\dot{\mathbf{k}}_G\approx 0$）、SRBD 把 $\mathbf{A}_G$ 的角动量部分冻结成常惯量 $\mathbf{I}_G$、Centroidal 则保留完整的 $\mathbf{A}_G(q)$。\textbf{“丢什么物理”这一谱系主线，在 CMM 视角下变成“近似 $\mathbf{A}_G$ 的哪几列／哪几块”} ——一句话统一了前面所有章节。
\end{insight}

\begin{table}[H]\centering\small
\caption{四档模型与质心动量矩阵 $\mathbf{A}_G$ 的关系——全章谱系的收束}
\label{tab:rm-cmm-bridge}
\begin{tabular}{lll}
\toprule
档位 & 对 $\mathbf{A}_G$／动量的处理 & 角动量建模 \\
\midrule
LIPM       & 只用线动量、$\dot{\mathbf{k}}_G\approx 0$ & 忽略 \\
SRBD       & 冻结角动量块为常惯量 $\mathbf{k}_G=\mathbf{I}_G^{\mathrm{const}}\boldsymbol{\omega}$ & 恒定 $\mathbf{I}_G$ \\
Centroidal & 保留完整 $\mathbf{k}_G=\mathbf{A}_G^{\mathrm{ang}}(q)\dot{q}$ & 精确 \\
全身动力学  & 完整 \cref{eq:rm-fullbody}（含关节级约束） & 精确 $+$ 关节约束 \\
\bottomrule
\end{tabular}
\end{table}

\begin{note}
$\mathbf{A}_G$ 的高效 $O(N)$ 计算（CCRBA）及复合刚体质心惯量 $\mathbf{I}_G=\mathbf{A}_G\mathbf{M}^{-1}\mathbf{A}_G^\top$ 归 \cref{ch:contact-mechanics}\cite{orin2013centroidal}；其空间力变换 ${}^{G}\!\mathbf{X}^*$ 的伴随来源见 \cref{ch:spatial-dynamics}。本节用 CMM 作“全身 $\leftrightarrow$ 简化”的概念桥，不重铺其算法；同一个 $\mathbf{A}_G$ 在 \cref{ch:contact-mechanics} 还承担“接触力传递”的用途，两处互引。
\end{note}

\begin{exercise}
\begin{enumerate}
  \item 解释为什么 \cref{eq:rm-fullbody} 的“基座 $6$ 行”不含关节力矩 $\boldsymbol{\tau}$，以及这如何对应 \cref{eq:rm-centroidal-dyn} 中“$\dot{\mathbf{h}}_G$ 只由接触力与重力决定”。
  \item 据 \cref{tab:rm-cmm-bridge}，把“SRBD 在跳跃时角动量误差大”重新表述为“对 $\mathbf{A}_G$ 的哪一块近似失效”。
  \item 本书取 $\mathbf{h}_G=[\mathbf{l}_G;\mathbf{k}_G]$（线在前），而 Orin/Featherstone 取 $[\mathbf{k};\mathbf{l}]$。写出二者之间的 $6\times 6$ 置换矩阵，并说明为何不影响 \cref{eq:rm-centroidal-dyn} 的物理内容。
\end{enumerate}
\end{exercise}

% ════════════════════════════════════════════════════════════════
\section{虚拟模型控制 VMC}
\label{sec:rm-vmc}

承上：谱系的下游接口（\cref{tab:rm-downstream}）总要把“期望力”落到关节。在 MPC/WBC 框架成熟之前，腿足控制曾用一种更轻、更直觉的力控方法——虚拟模型控制（VMC），至今仍是嵌入式平台与控制器调参的重要心智模型。它是谱系“向控制层落地”的一块独立拼图。

\paragraph{核心思想与力映射。}
VMC 用虚拟的机械元件（弹簧、阻尼器等）连接机器人作用点或连接作用点与环境，由虚拟元件产生\emph{虚拟力} $\mathbf{F}_{\mathrm{virtual}}$ 来“拉/推”机器人达成期望运动；再经 Jacobian 转置把工作空间虚拟力映成关节力矩：
\begin{equation}
  \boxed{\;\boldsymbol{\tau}=\mathbf{J}^\top(q)\,\mathbf{F}_{\mathrm{virtual}}\;}.
  \label{eq:rm-vmc}
\end{equation}
其优势是：无须逆动力学求解、只需正运动学与 Jacobian，计算高效、物理直觉清晰（“弹簧把身体往上拉”），适合嵌入式实时运行。

\begin{theorem}[VMC 力映射的虚功原理证明]
\label{thm:rm-vmc}
\cref{eq:rm-vmc} 不是近似，而是\textbf{精确的静力学等价关系}。
\end{theorem}
\begin{proof}
设关节空间广义力为 $\boldsymbol{\tau}$、工作空间力为 $\mathbf{F}$。虚功原理要求二者对同一虚位移做功相等：$\boldsymbol{\tau}^\top\delta q=\mathbf{F}^\top\delta\mathbf{x}$。由正运动学微分 $\delta\mathbf{x}=\mathbf{J}(q)\,\delta q$ 代入得 $\boldsymbol{\tau}^\top\delta q=\mathbf{F}^\top\mathbf{J}\,\delta q$。因 $\delta q$ 任意，故 $\boldsymbol{\tau}=\mathbf{J}^\top\mathbf{F}$。\qed
\end{proof}

\begin{insight}
\cref{thm:rm-vmc} 在\emph{静力学}下精确，但在动力学条件（加速度非零）下，VMC 的力映射\textbf{忽略了惯性力与科氏力}——这正是它在高动态场景精度不及 WBC 的根因。WBC（\cref{ch:control} 与 P7 的 WBC 章）通过求解完整逆动力学 $\boldsymbol{\tau}=\mathbf{M}\ddot{q}+\mathbf{h}-\mathbf{J}_c^\top\boldsymbol{\lambda}$ 精确计入惯性项。可以说，\textbf{VMC 是 WBC 在“非冗余、静力学主导”情形下的特化}：单腿非冗余系统里，VMC 的 Jacobian 转置力映射与 WBC 的接触力分配本质相同；VMC 无法处理的多任务优先级与复杂约束，正是 WBC 框架要解决的。误差随运动加速度增大而上升——静态站立时 VMC 与真值偏差极小，快速 trot 或跳跃时则因惯性项被忽略而显著（经验量级、依平台而变\rebuilt{VMC vs WBC 的具体误差百分比无原始出处，仅作趋势}）。
\end{insight}

\begin{note}
\cref{eq:rm-vmc} 的 Jacobian 转置力映射与 \cref{ch:spatial-dynamics} 的接触虚功对偶 $\boldsymbol{\tau}=\mathbf{J}_c^\top\boldsymbol{\lambda}$、以及操作空间力控（\cref{ch:control}）同源——转置约定见 \cref{ch:matderiv}。本节只补“VMC $=$ 把期望行为编码成虚拟力”的控制哲学，力映射的一般理论不在此展开。
\end{note}

\begin{exercise}
\begin{enumerate}
  \item 对二自由度单腿（连杆长 $L_1,L_2$、关节角 $q_1,q_2$），写出足端位置的正运动学与 Jacobian $\mathbf{J}(q)$，并据 \cref{eq:rm-vmc} 把足端虚拟力 $\mathbf{F}_{\mathrm{virtual}}=K_p(\mathbf{p}_d-\mathbf{p})+K_d(\dot{\mathbf{p}}_d-\dot{\mathbf{p}})$ 映成两个关节力矩。
  \item 解释为什么 \cref{thm:rm-vmc} 在静力学下精确、动力学下有误差，并指出误差项的物理量级随什么增大。
  \item 用一句话说明 VMC 与 WBC 的关系，以及 VMC 在现代 MPC$+$WBC 栈中被谁取代、其思想为何仍有价值。
\end{enumerate}
\end{exercise}

% ════════════════════════════════════════════════════════════════
\section{本章脉络与展望}
\label{sec:rm-outlook}

\paragraph{一张全景图。}
本章从维度灾难出发，沿“线性平动 $\to$ 一阶降阶 $\to$ 弹簧飞行相 $\to$ 转动 $\to$ 精确动量 $\to$ 全身”逐档走完降阶谱系，并用质心动量矩阵把它收束成一座桥。\cref{fig:rm-panorama} 是这张谱系全景图：横轴是\emph{维度／精度}递增，纵轴标注\emph{核心假设}与\emph{凸性}，CMM 作为贯穿全图的“全身 $\leftrightarrow$ 简化”投影。

\begin{figure}[H]
\centering
\begin{tikzpicture}[
  >=Stealth, font=\small, node distance=4mm,
  m/.style={draw, rounded corners, align=center, font=\scriptsize, inner sep=3pt, minimum height=11mm, text width=17mm},
  cvx/.style={m, fill=main!10, draw=main!60!black},
  ncvx/.style={m, fill=second!12, draw=second!60!black},
  full/.style={m, fill=third!12, draw=third!60!black},
  ln/.style={->, thick, gray!55!black}]
\node[cvx] (lip) {LIPM\\\scriptsize 恒高·点质量};
\node[cvx, right=of lip] (dcm) {DCM\\\scriptsize 一阶降阶};
\node[ncvx, right=of dcm] (slip) {SLIP\\\scriptsize 弹簧·飞行相};
\node[cvx, right=of slip] (srb) {SRBD\\\scriptsize 恒 $\mathbf{I}_G$·凸};
\node[cvx, right=of srb] (ll) {LL-SRBM\\\scriptsize 腿聚合·凸};
\node[ncvx, right=of ll] (cen) {Centroidal\\\scriptsize 精确 $\mathbf{A}_G$·NLP};
\node[full, right=of cen] (fb) {全身动力学\\\scriptsize 关节约束};
\foreach \a/\b in {lip/dcm, dcm/slip, slip/srb, srb/ll, ll/cen, cen/fb}{\draw[ln] (\a)--(\b);}
% axis
\draw[->, thick] (lip.south west)++(0,-6mm) -- ++(11.6,0) node[right, font=\scriptsize]{维度／精度递增};
% CMM bridge
\node[draw, dashed, rounded corners, fill=gray!8, font=\scriptsize, below=10mm of srb, text width=52mm, align=center]
  (cmm) {质心动量桥 $\mathbf{h}_G=\mathbf{A}_G(q)\dot{q}$\\（各档 $=$ 对 $\mathbf{A}_G$ 的不同近似）};
\draw[ln, dashed] (cmm.north) -- (srb.south);
\draw[ln, dashed] (cmm.north west) to[out=120,in=-60] (lip.south);
\draw[ln, dashed] (cmm.north east) to[out=60,in=-120] (fb.south);
\end{tikzpicture}
\caption{降阶模型谱系全景图：维度／精度自左向右递增，蓝＝凸（LIPM/DCM/SRBD/LL-SRBM）、橙＝非凸（SLIP/Centroidal）、绿＝全身；质心动量矩阵 $\mathbf{A}_G$ 是贯穿全谱的精确桥。}
\label{fig:rm-panorama}
\end{figure}

\paragraph{前沿：模型同伦迁移。}
谱系的一个新近方向是把“离散地选一档”变成“连续地在档间滑动”——\textbf{模型同伦迁移}（model homotopy transfer）。引入参数 $\alpha\in[0,1]$：$\alpha=0$ 时系统是纯 SRBD（全部质量在躯干），$\alpha=1$ 时是全身动力学（质量按真实分布）；沿光滑路径缓慢增大 $\alpha$，让优化器/策略从“已知的 SRBD 最优解”逐步追踪到“全身动力学的解”，比直接在高维全身模型上从零优化更稳定（思想源自同伦延续法）。这一方法在四足高动态任务（如 wall-jump、翻转）上已有报道\rebuilt{Kang et al., "Dynamic Policy Learning for Legged Robot with Simplified Model Pretraining and Model-Homotopy-Inspired Transfer", IEEE RA-L 2026（arXiv:2512.24698）；紧贴知识边界，标注待核}\cite{kang2026dynamic}。它把本章“离散谱系”升维成“连续模型空间”，是降阶模型与策略学习的交叉前沿。RL 作为“万能近似器”整体替代简化模型的范式，属强化学习运控部范畴，本章不展开。

\begin{insight}
回望全章，降阶建模的智慧不在“算得多精”，而在\textbf{知道自己丢了什么、并据任务决定丢得起丢不起}。LIPM 丢了转动却换来解析的预览控制与 DCM；SRBD 丢了惯量变化却换来毫秒级凸 QP；Centroidal 什么都不丢却付出非线性优化的代价。CMM 桥（\cref{eq:rm-cmm}）则保证：无论降到哪一档，都能精确地说清它是全身动力学的哪一个投影。这正是“地图”章的价值——它不替你走路，但让你知道每条路通向哪里、要付什么代价。
\end{insight}

\section*{常见误解汇总}
\addcontentsline{toc}{section}{常见误解汇总}
\begin{enumerate}
  \item \textbf{“简化模型可以直接驱动电机”}：错。降阶模型只活在规划层，输出的是期望质心力/接触力，须经 WBC 才能变成关节力矩（\cref{fig:rm-pyramid}、\cref{ch:control}）。
  \item \textbf{“ZMP 在支撑多边形内就一定不摔倒”}：错。这是\emph{必要非充分}条件，只管住“绕边界翻倒”，不防滑动、力矩饱和等（\cref{thm:rm-zmp-stable} 后的陷阱）。
  \item \textbf{“收敛分量就是 VRP”}：错。收敛分量 $\zeta=x-\dot{x}/\omega_n$（\cref{eq:rm-ccm}）是\emph{衰减}模态；VRP 是\emph{排斥}点（编码总力），二者不可混（订正自部分文献的误植）。
  \item \textbf{“变高度行走可套用恒高 DCM”}：错。$h$ 变化时须用变高度 DCM 或 VHIP（\cref{ch:legged-biped}、Caron 2019\cite{caron2019capturability}）；不存在所谓“VRPM”这种标准模型。
  \item \textbf{“SLIP 的弹簧是物理弹簧”}：错。$k$ 是由关节阻抗/力控产生的\emph{等效刚度}（\cref{sec:rm-slip} 陷阱）。
  \item \textbf{“SRBD 与 Centroidal 在四足 trot 上差别很大”}：未必。CoM 位置预测差别小（腿质量占比小），但\emph{角动量}预测差别显著——是否升档取决于任务是否在乎角动量（\cref{tab:rm-infoloss}）。
  \item \textbf{“Centroidal 最精确所以总该用它”}：错。它是非线性规划、算力贵；$e_{\mathrm{model}}<5\%$ 的低速行走用 SRBD 更划算（\cref{eq:rm-emodel} 的升降档规则）。
\end{enumerate}

\section*{小结}
\addcontentsline{toc}{section}{小结}
本章把腿足规划层用到的全部降阶模型组织成一张\emph{谱系}：
\begin{itemize}
  \item \textbf{维度灾难}（\cref{sec:rm-curse}）：全身 MPC 维度爆炸（\cref{eq:rm-decision-count}），故分“规划层降阶 MPC $+$ 控制层 WBC”两级；谱系总表 \cref{tab:rm-landscape} 是全章主线锚。
  \item \textbf{LIPM}（\cref{sec:rm-lipm}）：$\ddot{x}=\omega_n^2(x-p)$（\cref{eq:rm-lipm}）是“二阶不稳定 ODE 范式”，本征不稳定、含双曲离散解（\cref{eq:rm-lipm-disc}）；是 $\omega_n$/ZMP/DCM/SLIP 的公共地基。
  \item \textbf{CoM/CoP/ZMP}（\cref{sec:rm-zmp}）：平地 ZMP$=$CoP（\cref{thm:rm-zmp-cop}）、含角动量项的完整 ZMP 公式（\cref{eq:rm-zmp-full}）退化为 LIPM；判据必要非充分。
  \item \textbf{ZMP 预览}（\cref{sec:rm-preview}）：Cart-Table $+$ LQR $+$ preview，预览增益指数衰减（\cref{eq:rm-preview-gain}），窗口 $\approx1.5$\,s 足够（Kajita 2003 / HRP-2）。
  \item \textbf{DCM}（\cref{sec:rm-dcm}）：$\xi=x+\dot{x}/\omega_n$ 把二阶降一阶（\cref{eq:rm-dcm-dyn}）、捕获点 $=$ 当前 DCM；本质是“把控制资源集中于不稳定模态”。
  \item \textbf{SLIP}（\cref{sec:rm-slip}）：弹簧腿 $+$ 飞行相补回垂直能量，极坐标方程无闭式解（\cref{eq:rm-slip}），Poincaré 判稳、被动自稳定；Raibert 落脚律（\cref{eq:rm-raibert}）$\approx$ 一阶捕获点。
  \item \textbf{SRBD/Centroidal}（\cref{sec:rm-srbd-centroidal}）：补回转动；SRBD 恒 $\mathbf{I}_G$ 故凸（推导见 \cref{ch:quad_mpc}），Centroidal 精确 $\mathbf{A}_G(q)$ 故 NLP（Orin 2013 洞察）。
  \item \textbf{选型理论}（\cref{sec:rm-selection}）：信息丢失表 \cref{tab:rm-infoloss} 是灵魂，LL-SRBM 是凸 QP 中间档（\cref{eq:rm-llsrbm}），$e_{\mathrm{model}}$ 给升降档定量依据。
  \item \textbf{CMM 桥}（\cref{sec:rm-bridge}）：$\mathbf{h}_G=\mathbf{A}_G(q)\dot{q}$（\cref{eq:rm-cmm}）把每档解释成“对 $\mathbf{A}_G$ 的某种近似”（\cref{tab:rm-cmm-bridge}），收束全谱。
  \item \textbf{VMC}（\cref{sec:rm-vmc}）：$\boldsymbol{\tau}=\mathbf{J}^\top\mathbf{F}$ 虚功精确（\cref{thm:rm-vmc}），动力学下忽略惯性、是 WBC 的特化。
\end{itemize}

\section*{延伸阅读}
\addcontentsline{toc}{section}{延伸阅读}
\begin{itemize}
  \item \textbf{LIPM 与 ZMP 预览}：Kajita \& Tani 1991\cite{kajita1991study}（LIPM 首用）；Kajita et al.\ 2001\cite{kajita2001lipm}（3D-LIPM 与预览控制谱系）；ZMP 史与判据见 Vukobratović \& Borovac 2004\cite{vukobratovic2004zero}。
  \item \textbf{DCM 与可捕获性}：Pratt et al.\ 2006\cite{pratt2006capture}（捕获点）；Koolen et al.\ 2012\cite{koolen2012capturability}（可捕获性分析）；Englsberger et al.\ 2015\cite{englsberger2015dcm}（3D-DCM/eCMP/VRP）；变高度见 Caron 2019\cite{caron2019capturability}（VHIP）。DCM 工程化深挖见 \cref{ch:legged-biped}。
  \item \textbf{SLIP 与 Raibert}：Raibert 1986\cite{raibert1986legged}（三分解耦与落脚启发式）；Geyer, Seyfarth \& Blickhan 2006\cite{geyer2006compliant}（弹簧–质量统一行走奔跑、被动自稳定）。
  \item \textbf{SRBD 与 Centroidal}：Di Carlo et al.\ 2018\cite{dicarlo2018cheetah}（MIT Cheetah 3 凸 MPC，全栈见 \cref{ch:quad_mpc}）；Orin, Goswami \& Lee 2013\cite{orin2013centroidal}（质心动力学与 CMM）；Wensing \& Orin 2016\cite{wensing2016improved}（改进的 $O(N)$ 质心动力学计算）。
  \item \textbf{综述与选型}：Wensing et al.\ 2024\cite{wensing2024legged}（动态腿足机器人优化控制综述，覆盖 LIPM/SRBD/Centroidal 全谱与精度–算力权衡）；前沿模型同伦迁移见 Kang et al.\ 2026\cite{kang2026dynamic}。
\end{itemize}

\section*{本章与后续章节的关系}
\addcontentsline{toc}{section}{本章与后续章节的关系}
\begin{itemize}
  \item $\to$ \cref{ch:quad_mpc}（单刚体 MPC）：SRBD 完整推导、三大凸化、$13$ 维状态、堆叠 QP、Di Carlo 工程——本章 \cref{subsec:rm-srbd} 的谱系定位在此落到全栈实现。
  \item $\to$ \cref{ch:legged-biped}（双足点足控制）：DCM 步态反向递推、反馈律、ALIP 横向落脚——本章 \cref{sec:rm-dcm} 的一阶降阶洞察在此工程化。
  \item $\to$ \cref{ch:spatial-dynamics}（空间向量/RNEA）：全身动力学 \cref{eq:rm-fullbody}、选择矩阵块分解、接触 Jacobian——本章 \cref{sec:rm-bridge} 只引用、不重证。
  \item $\to$ \cref{ch:contact-mechanics}（接触力学）：CCRBA 的 $O(N)$ 算法、摩擦锥、GIWC——本章 CMM 桥 \cref{eq:rm-cmm} 取其算法。
  \item $\to$ \cref{ch:control}（控制导论）：WBC/HQP、操作空间力控——本章 \cref{sec:rm-vmc} 的 VMC 是其特化。
  \item $\to$ \cref{ch:lqr-lqg-riccati}（LQR/Riccati）：本章 \cref{sec:rm-preview} 的预览控制是“LQR $+$ preview”的腿足应用；\cref{ch:nlopt}/\cref{ch:ddp-ilqr} 提供 Centroidal NLP 的求解理论。
\end{itemize}
